% ===== PRÉAMBULE CANONIQUE — à coller VERBATIM en tête de CHAQUE .tex =====
\documentclass[11pt,a4paper]{article}
\usepackage[utf8]{inputenc}\usepackage[T1]{fontenc}\usepackage{lmodern}
\usepackage[french]{babel}
\usepackage{amsmath,amssymb,mathtools}\usepackage{icomma}
\usepackage{siunitx}\sisetup{locale=FR}  % \qty{}{}, \si{}, \ang{}
\usepackage{xcolor}\usepackage{colortbl}
\usepackage{tikz}\usepackage{pgfplots}\pgfplotsset{compat=1.18}
\usepackage[siunitx,europeanresistors]{circuitikz} % circuits (élec.)
\usepackage[most]{tcolorbox}\usepackage{enumitem}\usepackage{array,booktabs}
\usepackage[a4paper,margin=2.2cm]{geometry}
\usetikzlibrary{arrows.meta,calc,angles,quotes,positioning,patterns,%
  backgrounds,shapes.geometric,decorations.pathmorphing,decorations.markings,intersections}
\definecolor{primary}{RGB}{222,49,99}\definecolor{primarydark}{RGB}{150,22,55}
\definecolor{defcol}{RGB}{0,114,178}\definecolor{methcol}{RGB}{52,92,148}
\definecolor{excol}{RGB}{0,135,90}\definecolor{attncol}{RGB}{200,40,55}
\tcbset{boxbase/.style={enhanced,breakable,boxrule=0.9pt,arc=2.5pt,
  left=3mm,right=3mm,top=2.3mm,bottom=2.3mm,fonttitle=\bfseries,coltitle=white}}
% --- boîtes TUTORIEL (noms EXACTS, ne pas renommer) ---
\newtcolorbox{cle}[1][]{boxbase,colback=defcol!6,colframe=defcol,
  title={\textsc{À retenir}\ifx\relax#1\relax\relax\else\textmd{ : \itshape #1}\fi}}
\newtcolorbox{methode}[1][]{boxbase,colback=methcol!7,colframe=methcol,sharp corners,
  title={\textsc{Méthode}\ifx\relax#1\relax\relax\else\textmd{ --- \itshape #1}\fi}}
\newtcolorbox{exemplebox}[1][]{boxbase,colback=excol!5,colframe=excol,
  title={\textsc{Exemple}\ifx\relax#1\relax\relax\else\textmd{ : \itshape #1}\fi}}
\newtcolorbox{piege}[1][]{boxbase,colback=attncol!6,colframe=attncol,
  title={\textsc{Piège}\ifx\relax#1\relax\relax\else\textmd{ : \itshape #1}\fi}}
\newtcolorbox{remarque}[1][]{boxbase,colback=black!4,colframe=black!45,
  title={\textsc{Remarque}\ifx\relax#1\relax\relax\else\textmd{ : \itshape #1}\fi}}
% --- boîtes EXERCICE / CORRIGÉ (noms EXACTS, auto-comptées) ---
\newtcolorbox[auto counter]{exo}[1][]{boxbase,colback=primary!4,colframe=primary,
  title={\textsc{Exercice}~\thetcbcounter\ifx\relax#1\relax\relax\else\textmd{ --- \itshape #1}\fi}}
\newtcolorbox[auto counter]{corrige}[1][]{boxbase,colback=excol!4,colframe=excol,
  title={\textsc{Corrigé}~\thetcbcounter\ifx\relax#1\relax\relax\else\textmd{ --- \itshape #1}\fi}}
\newcommand{\vv}[1]{\vec{#1}}\newcommand{\ds}{\displaystyle}
\newcommand{\R}{\mathbb{R}}\newcommand{\N}{\mathbb{N}}\newcommand{\Z}{\mathbb{Z}}
% =========================================================================
\begin{document}
\sisetup{per-mode=symbol}

\begin{exo}[chute non libre : l'accélération]
Un parachutiste (masse $m=\qty{80}{\kilogram}$) tombe. À un instant donné, la résistance de l'air
vaut $F_f=\qty{300}{\newton}$ ; on prend $g=\qty{10}{\metre\per\second\squared}$.
\begin{enumerate}[label=\alph*)]
  \item Écris le bilan des forces dans le sens de la chute.
  \item Calcule l'accélération à cet instant.
\end{enumerate}
\end{exo}

\begin{exo}[retrouver la résistance de l'air]
Une bille de masse $m=\qty{2}{\kilogram}$ tombe dans l'air avec une accélération mesurée
$a=\qty{8}{\metre\per\second\squared}$ ($g=\qty{10}{\metre\per\second\squared}$).
Quelle est l'intensité de la résistance de l'air $F_f$ à cet instant ?
\end{exo}

\begin{exo}[la pomme de pin à vitesse constante]
Une pomme de pin de masse $m=\qty{50}{\gram}$ tombe d'un arbre. Après quelques secondes, elle atteint
une \emph{vitesse constante} de \qty{3}{\metre\per\second}. On prend $g=\qty{10}{\metre\per\second\squared}$.
\begin{enumerate}[label=\alph*)]
  \item Que vaut son accélération quand la vitesse est constante ? Qu'en déduit-on sur les forces ?
  \item Calcule alors l'intensité de la force de frottement de l'air.
\end{enumerate}
\end{exo}

\begin{exo}[pourquoi $a<g$ ?]
Un objet de masse $m=\qty{2}{\kilogram}$ tombe dans l'air. Son poids vaut $G=\qty{20}{\newton}$ et la
résistance de l'air $F_f=\qty{4}{\newton}$.
\begin{enumerate}[label=\alph*)]
  \item Calcule l'accélération de l'objet.
  \item Compare-la à $g=\qty{10}{\metre\per\second\squared}$ et explique pourquoi elle est plus petite.
\end{enumerate}
\end{exo}

\begin{exo}[deux masses lancées sur la Lune]
Sur la Lune (pas d'atmosphère, $g=\qty{1,6}{\metre\per\second\squared}$), on lance verticalement vers le
haut, \emph{en même temps} et à la \emph{même vitesse}, deux masses $m_1=\qty{3}{\kilogram}$ et
$m_2=\qty{6}{\kilogram}$.
\begin{enumerate}[label=\alph*)]
  \item Les deux masses ont-elles la même accélération pendant le vol ? Justifie.
  \item Laquelle monte le plus haut ?
\end{enumerate}
\end{exo}
\end{document}
